%% kapa: exam=1 round=19 year=2008 period=1 form=A subject=civil_law subject_ko=민법 seq=31 number=31
%% source: Q-Net 감정평가사 기출문제 (cst00309) artlSeq=1001962 file=extracted/1008/1001962/2008년 제19회 감정평가사 1차시험 기출문제/1교시.hwp
%% conversion: kordoc 4.24.51 -> tools/gen1.py (auto); review=proofread-sonnet; flags=none
\begin{question}{31}
등기부취득시효에 관한 설명으로 틀린 것을 모두 고른 것은? (다툼이 있는 경우에는 판례에 의함)
\begin{framedtext}
ㄱ. 점유자가 소유자로 등기되어 있어야 하는데, 이때의 등기는 적법․유효한 등기이어야 한다. \par
ㄴ. 등기기간은 전 소유자의 등기기간까지 포함하여 10년이면 된다. \par
ㄷ. 공유자의 1인이 공유부동산 중 특정부분만을 점유하여 왔다면 그 특정부분에 대한 공유지분의 범위 내에서만 등기부 취득시효할 수 있다. \par
ㄹ. 소유권보존등기가 이중으로 경료되어 뒤에 된 소유권보존등기가 무효로 되는 때에도, 뒤에 된 소유권보존등기를 기초로 한 등기부취득시효의 완성을 주장할 수 있다. \par
ㅁ. 선의․무과실은 등기에 관한 것이 아니라 점유취득에 관한 것이며, 무과실에 관한 증명책임은 시효취득을 주장하는 자에게 있다.
\end{framedtext}
\begin{choices}
  \item ㄱ. ㄴ
  \item ㄱ, ㄹ
  \item ㄷ, ㅁ
  \item ㄱ, ㄴ, ㄹ
  \item ㄷ, ㄹ, ㅁ
\end{choices}
\end{question}
\answer{2}
